\section{Induction: Infinite Claims, Finite Proofs}
\chapterline

Imagine standing before an infinite line of dominoes stretching past the horizon:
\[
D_1, \; D_2, \; D_3, \; D_4, \; \dots, \; D_n, \; \dots
\]
You wish to ensure that every single domino will eventually topple. 

You could walk down the line and knock each domino over by hand. But because the line is infinite, you would never finish. No matter how many millions of dominoes you tip over, an endless expanse remains untouched.

Instead, you inspect the mechanics of the line and notice two local facts:
\begin{enumerate}[label=\arabic*.]
\item \textbf{The Initial Push:} You physically tip over the very first domino, $D_1$.
\item \textbf{The Chain Reaction:} The dominoes are spaced closely enough that whenever any arbitrary domino $D_k$ falls, its momentum is guaranteed to strike and topple the next domino, $D_{k+1}$.
\end{enumerate}

With those two guarantees in place, every domino in the infinite chain falls: domino $1$ topples $2$, domino $2$ topples $3$, and the pattern continues for every positive integer.

This is the principle of \textbf{mathematical induction}. It is our primary tool for proving an infinite family of statements using a finite chain of deductions.

\begin{historybox}
The formalization of mathematical induction was achieved by Blaise Pascal in his 1654 \emph{Trait\'e du triangle arithm\'etique}. Pascal was collaborating with Pierre de Fermat on the ``problem of points''---how to fairly divide the stakes in an unfinished game of chance. To calculate probabilities, Pascal needed rigorous formulas for binomial coefficients across all rows of his triangle, but he had no way to verify an infinite table one row at a time. He formulated the principle of induction: verify the first row, prove that each row forces the truth of the next, and the entire table is established.

Traces of inductive reasoning had appeared earlier in the work of the 14th-century mathematician Levi ben Gershon (Gersonides) and the Sicilian mathematician Francesco Maurolico (1575), who proved that the sum of the first $n$ odd integers equals $n^2$. But Pascal made the principle explicit and general. 

Two centuries later, Giuseppe Peano (\emph{Arithmetices principia}, 1889) recognized that induction is not just a proof technique; it is the fifth Peano axiom defining what the natural numbers $\N$ are. The natural numbers are precisely the set generated from $1$ by repeatedly applying the successor operation. Induction works because there are no natural numbers outside the sequence generated by repeated succession.
\end{historybox}

\begin{center}
\begin{tikzpicture}[scale=0.92, >=Latex]
  % Stage box
  \draw[rounded corners=4pt, fill=gray!3, draw=gray!25] (-0.6, -1.15) rectangle (12.5, 3.4);
  \node[anchor=west, font=\small\bfseries] at (-0.3, 3.0) {Mathematical Induction as an Infinite Domino Cascade};
  
  % Ground line
  \draw[thick, gray!50] (0.0, 0.0) -- (11.6, 0.0);
  
  % Domino 1: Tipping over (Base case)
  \begin{scope}[shift={(1.0, 0.0)}, rotate around={-25:(0,0)}]
    \draw[thick, sourceblue, fill=sourceblue!25, rounded corners=1pt] (0, 0) rectangle (0.28, 1.8);
    \node[sourceblue!90!black, font=\scriptsize\bfseries] at (0.14, 0.9) {$1$};
  \end{scope}
  \draw[->, very thick, trapred] (0.0, 1.5) -- node[above, font=\scriptsize\bfseries, text=trapred!90!black] {Base Push} (0.8, 1.2);
  \node[below, font=\scriptsize, align=center, text=sourceblue!90!black] at (1.0, -0.2) {Base Case\\$P(1)$ verified};

  % Domino 2: Tipping
  \begin{scope}[shift={(2.2, 0.0)}, rotate around={-20:(0,0)}]
    \draw[thick, sourceblue, fill=sourceblue!20, rounded corners=1pt] (0, 0) rectangle (0.28, 1.8);
    \node[sourceblue!90!black, font=\scriptsize\bfseries] at (0.14, 0.9) {$2$};
  \end{scope}
  \node[below, font=\scriptsize, text=gray!70] at (2.2, -0.2) {$P(2)$};

  % Domino 3: Tipping
  \begin{scope}[shift={(3.4, 0.0)}, rotate around={-15:(0,0)}]
    \draw[thick, sourceblue, fill=sourceblue!15, rounded corners=1pt] (0, 0) rectangle (0.28, 1.8);
    \node[sourceblue!90!black, font=\scriptsize\bfseries] at (0.14, 0.9) {$3$};
  \end{scope}
  \node[below, font=\scriptsize, text=gray!70] at (3.4, -0.2) {$P(3)$};

  % Dots
  \node[font=\large\bfseries, text=gray!50] at (4.8, 0.9) {$\dots$};

  % Domino k (Hypothesis)
  \begin{scope}[shift={(5.8, 0.0)}, rotate around={-15:(0,0)}]
    \draw[thick, quotientorange, fill=quotientorange!25, rounded corners=1pt] (0, 0) rectangle (0.28, 1.8);
    \node[quotientorange!90!black, font=\scriptsize\bfseries] at (0.14, 0.9) {$k$};
  \end{scope}
  \node[below, font=\scriptsize, align=center, text=quotientorange!90!black] at (5.8, -0.2) {Assume $P(k)$\\(Hypothesis)};

  % Domino k+1 (Target)
  \begin{scope}[shift={(8.0, 0.0)}]
    \draw[thick, targetgreen, fill=targetgreen!25, rounded corners=1pt] (0, 0) rectangle (0.38, 1.8);
    \node[targetgreen!90!black, font=\tiny\bfseries] at (0.19, 0.9) {$k\!+\!1$};
  \end{scope}
  \node[below, font=\scriptsize, align=center, text=targetgreen!90!black] at (8.19, -0.2) {Deduce $P(k\!+\!1)$\\(Conclusion)};

  % Inductive Bridge Arc
  \draw[->, very thick, structurepurple] (6.1, 2.0) to[bend left=35] node[above, font=\scriptsize\bfseries, text=structurepurple!90!black] {Inductive Step: $P(k) \implies P(k+1)$} (8.15, 2.0);

  % Domino k+2 and beyond
  \begin{scope}[shift={(9.6, 0.0)}]
    \draw[thick, gray!50, fill=gray!10, rounded corners=1pt] (0, 0) rectangle (0.28, 1.8);
  \end{scope}
  \node[font=\large\bfseries, text=gray!40] at (10.2, 0.9) {$\dots$};
  \draw[->, very thick, gray!40] (10.7, 0.9) -- (11.3, 0.9) node[right, font=\tiny, text=gray!60] {to $\infty$};
\end{tikzpicture}\\[0.4em]
{\small\textbf{Figure 17.1:} Mathematical induction as an infinite cascade. An initial manual push establishes the base case $P(1)$. The conditional guarantee $P(k) \implies P(k+1)$ acts like the close spacing between dominoes: whenever any domino $k$ falls, it knocks over $k+1$. Together, these two finite proofs guarantee that every domino in the infinite sequence falls.}
\end{center}

\subsection{The Formal Principle of Mathematical Induction}

To make the domino intuition mathematically rigorous, we replace dominoes with mathematical statements.

Recall from Chapter~2 that a \textbf{predicate} is an open statement containing a variable slot, such as $P(n): \text{``}2^n \ge n+1\text{''}$. By itself, a predicate is neither true nor false; it converts into a definite assertion with a truth value only when a specific number is plugged into the slot.

\begin{theorem}[Principle of Mathematical Induction]
Let $P(n)$ be a predicate parameterized by integers $n \ge n_0$. Suppose that:
\begin{enumerate}[label=(\roman*)]
\item \textbf{Base Case:} $P(n_0)$ is true.
\item \textbf{Inductive Step:} For every integer $k \ge n_0$, if $P(k)$ is true, then $P(k+1)$ is true:
\[
(\forall k \ge n_0)(P(k) \implies P(k+1)).
\]
\end{enumerate}
Then $P(n)$ is true for all integers $n \ge n_0$.
\end{theorem}

\begin{walkbox}
\textbf{Mental Check:} Suppose you verify the base case at $n_0 = 4$, and you prove that $P(k) \implies P(k+1)$ for all $k \ge 4$. Does this theorem guarantee that $P(2)$ is true?
\emph{Answer:} No. The domino cascade only travels forward from the base case. Dominoes prior to the starting push ($n < n_0$) remain untouched and unverified.
\end{walkbox}

\begin{intuitionbox}
A common question when first learning induction is: ``In the inductive step, aren't we assuming what we set out to prove?''

No: you are not asserting that $P(n)$ is universally true. You are proving an implication: if $P(k)$ holds, then $P(k+1)$ holds.

In the inductive step, you assume $P(k)$ hypothetically for an arbitrary $k \ge n_0$, and use that assumption to deduce $P(k+1)$. You are proving the validity of the conditional statement $P(k) \implies P(k+1)$. The base case then establishes the starting truth $P(n_0)$, allowing the truth of the statement to propagate to all subsequent integers.
\end{intuitionbox}

\subsection{A Worked Example: Gauss's Triangular Sum}

Consider the classical formula for the sum of the first $n$ positive integers:
\[
1 + 2 + 3 + \dots + n = \frac{n(n+1)}{2}.
\]
Before writing down any algebra, pause to see the visual geometry. If you arrange dots in a triangular staircase of height $n$, you can fit two such staircases together to form an $n \times (n+1)$ rectangle containing $n(n+1)$ dots. Half of that rectangle gives the total count (Figure~17.2).

\begin{center}
\begin{tikzpicture}[>=stealth, font=\small]
  % Card background container
  \draw[fill=white, draw=gray!25, rounded corners=6pt] (-0.8, -0.9) rectangle (11.5, 4.8);
  \draw[dashed, gray!30] (5.2, -0.7) -- (5.2, 4.5);

  % Panel 1: Single Staircase
  \node[font=\small\bfseries, sourceblue!85!black] at (2.2, 4.35) {1. Single Staircase ($S_5$)};

  % Panel 1 Dots (Blue)
  \foreach \c in {1} { \fill[sourceblue!85!black] ({0.7 + (\c-1)*0.65}, 3.3) circle (2.8pt); }
  \foreach \c in {1,2} { \fill[sourceblue!85!black] ({0.7 + (\c-1)*0.65}, 2.7) circle (2.8pt); }
  \foreach \c in {1,2,3} { \fill[sourceblue!85!black] ({0.7 + (\c-1)*0.65}, 2.1) circle (2.8pt); }
  \foreach \c in {1,2,3,4} { \fill[sourceblue!85!black] ({0.7 + (\c-1)*0.65}, 1.5) circle (2.8pt); }
  \foreach \c in {1,2,3,4,5} { \fill[sourceblue!85!black] ({0.7 + (\c-1)*0.65}, 0.9) circle (2.8pt); }

  % Panel 1 Row counts
  \node[right, font=\scriptsize, text=sourceblue!85!black] at (1.05, 3.3) {$1$ dot};
  \node[right, font=\scriptsize, text=sourceblue!85!black] at (1.70, 2.7) {$2$ dots};
  \node[right, font=\scriptsize, text=sourceblue!85!black] at (2.35, 2.1) {$3$ dots};
  \node[right, font=\scriptsize, text=sourceblue!85!black] at (3.00, 1.5) {$4$ dots};
  \node[right, font=\scriptsize, text=sourceblue!85!black] at (3.65, 0.9) {$5$ dots};

  % Panel 1 Dimension bar
  \draw[|<->|, >=stealth, gray!60, semithick] (0.1, 0.9) -- (0.1, 3.3) node[midway, left=2pt, font=\scriptsize, text=gray!80] {$n = 5$};

  % Panel 1 Summary
  \node[font=\footnotesize, align=center, text=sourceblue!90!black] at (2.2, 0.1) {Sum: $S_n = 1 + 2 + 3 + 4 + 5 = 15$};

  % Panel 2: Interlocked Rectangle
  \node[font=\small\bfseries, structurepurple!85!black] at (8.4, 4.35) {2. Interlocked Rectangle ($2 \cdot S_5$)};

  % Panel 2 Blue Dots (First Staircase)
  \foreach \c in {1} { \fill[sourceblue!85!black] ({5.8 + (\c-1)*0.65}, 3.3) circle (2.8pt); }
  \foreach \c in {1,2} { \fill[sourceblue!85!black] ({5.8 + (\c-1)*0.65}, 2.7) circle (2.8pt); }
  \foreach \c in {1,2,3} { \fill[sourceblue!85!black] ({5.8 + (\c-1)*0.65}, 2.1) circle (2.8pt); }
  \foreach \c in {1,2,3,4} { \fill[sourceblue!85!black] ({5.8 + (\c-1)*0.65}, 1.5) circle (2.8pt); }
  \foreach \c in {1,2,3,4,5} { \fill[sourceblue!85!black] ({5.8 + (\c-1)*0.65}, 0.9) circle (2.8pt); }

  % Panel 2 Orange Dots (Inverted Second Staircase)
  \foreach \c in {2,3,4,5,6} { \fill[quotientorange!90!black] ({5.8 + (\c-1)*0.65}, 3.3) circle (2.8pt); }
  \foreach \c in {3,4,5,6} { \fill[quotientorange!90!black] ({5.8 + (\c-1)*0.65}, 2.7) circle (2.8pt); }
  \foreach \c in {4,5,6} { \fill[quotientorange!90!black] ({5.8 + (\c-1)*0.65}, 2.1) circle (2.8pt); }
  \foreach \c in {5,6} { \fill[quotientorange!90!black] ({5.8 + (\c-1)*0.65}, 1.5) circle (2.8pt); }
  \foreach \c in {6} { \fill[quotientorange!90!black] ({5.8 + (\c-1)*0.65}, 0.9) circle (2.8pt); }

  % Panel 2 Dimension bars
  \draw[|<->|, >=stealth, gray!60, semithick] (5.8, 3.75) -- (9.05, 3.75) node[midway, above=2pt, font=\scriptsize, text=gray!80] {$n + 1 = 6$ columns};
  \draw[|<->|, >=stealth, gray!60, semithick] (9.65, 0.9) -- (9.65, 3.3) node[midway, right=2pt, font=\scriptsize, text=gray!80] {$n = 5$ rows};

  % Panel 2 Summary
  \node[font=\footnotesize, align=center, text=structurepurple!90!black] at (8.4, 0.1) {Total dots: $n(n+1) = 5 \times 6 = 30$\\[3pt]\textbf{One staircase:} $S_n = \dfrac{n(n+1)}{2} = 15$};

  % Legend below
  \node[font=\scriptsize, text=gray!70] at (5.3, -0.6) {{\color{sourceblue!85!black}$\bullet$} First copy $S_n$ \qquad {\color{quotientorange!90!black}$\bullet$} Second (inverted) copy $S_n$};
\end{tikzpicture}\\[0.4em]
{\small\textbf{Figure 17.2:} Visual proof of Gauss's triangular sum $\sum_{i=1}^n i = \frac{n(n+1)}{2}$. \emph{Left:} Dots arranged in a triangular staircase of height $n=5$, where row $i$ contains $i$ dots. \emph{Right:} Two congruent copies of the staircase (blue and orange) interlock to form an $n \times (n+1)$ rectangular grid. The rectangle contains $n(n+1)$ dots, so each staircase contains exactly half: $\frac{n(n+1)}{2}$.}
\end{center}

Now let us prove this identity rigorously for every integer $n \ge 1$ using mathematical induction.

\begin{theorem}
For all integers $n \ge 1$,
\[
\sum_{i=1}^n i = \frac{n(n+1)}{2}.
\]
\end{theorem}

\begin{proof}
Let $P(n)$ be the statement that $\sum_{i=1}^n i = \frac{n(n+1)}{2}$.

\textbf{Base Case ($n = 1$):}
For $n = 1$, the left-hand side is the single term $1$.
The right-hand side is
\[
\frac{1(1 + 1)}{2} = \frac{1(2)}{2} = 1.
\]
Both sides evaluate to $1$, so the base case $P(1)$ is verified.

\textbf{Inductive Step:}
Let $k \ge 1$ be an arbitrary integer. Assume that $P(k)$ is true, meaning
\[
1 + 2 + \dots + k = \frac{k(k+1)}{2}. \tag{Induction Hypothesis}
\]
We must prove that $P(k+1)$ holds, which means establishing:
\[
1 + 2 + \dots + k + (k+1) = \frac{(k+1)((k+1)+1)}{2} = \frac{(k+1)(k+2)}{2}.
\]

We begin with the sum of the first $k+1$ terms and isolate the sum of the first $k$ terms:
\[
\underbrace{1 + 2 + \dots + k}_{= \frac{k(k+1)}{2} \text{ by hypothesis}} + (k+1).
\]
Substituting our induction hypothesis:
\[
1 + 2 + \dots + k + (k+1) = \frac{k(k+1)}{2} + (k+1).
\]
Rather than expanding everything into quadratic polynomials, we notice the common factor of $(k+1)$ and factor it out cleanly:
\begin{align*}
\frac{k(k+1)}{2} + (k+1) 
&= (k+1)\left(\frac{k}{2} + 1\right) \\
&= (k+1)\left(\frac{k+2}{2}\right) \\
&= \frac{(k+1)(k+2)}{2}.
\end{align*}
This is precisely the formula demanded by $P(k+1)$.

Having verified the base case and established that $P(k) \implies P(k+1)$ for all $k \ge 1$, the Principle of Mathematical Induction implies that the formula holds for all integers $n \ge 1$.
\end{proof}

\begin{walkbox}
\textbf{Mental Check:} In the inductive step above, why did we peel off the last term $(k+1)$ before doing any algebra?
\emph{Answer:} Because our sole point of leverage is the induction hypothesis, which only speaks about sums of length $k$. Isolating the first $k$ terms creates the exact slot where our hypothesis can be legally applied.
\end{walkbox}

\subsection{Induction With Inequalities}

Induction is not limited to algebraic equalities; it frequently establishes growth rates and bounds in analysis.

When proving an equality by induction, both sides must match exactly. With inequalities, the induction step often yields a quantity strictly larger than needed. The goal is to bound terms in the correct direction rather than maintaining equality.

\begin{theorem}
For all integers $n \ge 1$, $2^n \ge n + 1$.
\end{theorem}

\begin{proof}
Let $P(n)$ be the inequality $2^n \ge n + 1$.

\textbf{Base Case ($n = 1$):} For $n = 1$, the left side is $2^1 = 2$ and the right side is $1 + 1 = 2$. Since $2 \ge 2$, the base case $P(1)$ holds.

\textbf{Inductive Step:} Let $k \ge 1$ be an arbitrary integer. Assume that $P(k)$ is true:
\[
2^k \ge k + 1. \tag{Induction Hypothesis}
\]
We must prove $P(k+1)$: that $2^{k+1} \ge (k+1) + 1 = k + 2$.

We start with the left side of our target inequality, expressing $2^{k+1}$ in terms of $2^k$:
\[
2^{k+1} = 2 \cdot 2^k.
\]
Because $2 > 0$, multiplying both sides of our induction hypothesis by $2$ preserves the inequality:
\[
2^{k+1} = 2 \cdot 2^k \ge 2(k + 1) = 2k + 2.
\]
Now we split $2k + 2$ to reveal our desired target $k + 2$:
\[
2k + 2 = k + k + 2.
\]
Because $k \ge 1$, the extra term $k$ is at least $1$. Therefore:
\[
2^{k+1} \ge k + k + 2 \ge 1 + k + 2 = k + 3 > k + 2.
\]
Since $2^{k+1} > k + 2$, it certainly follows that $2^{k+1} \ge k + 2$. This completes the inductive step.

By the Principle of Mathematical Induction, $2^n \ge n + 1$ for all integers $n \ge 1$.
\end{proof}

\subsection{Induction on Divisibility Statements}

A third common application of induction is divisibility in number theory.

Recall the definition from Chapter~6: for integers $a$ and $b$, we say $a \mid b$ ($a$ divides $b$) if there exists an integer $m$ such that $b = a \cdot m$.

When proving by induction that an algebraic expression is divisible by $d$, a systematic strategy is:
\begin{enumerate}[label=\arabic*.]
\item \textbf{State the Hypothesis:} Write the $k$-th expression as $d \cdot m$ for some integer $m$.
\item \textbf{Expand the Target:} Write out the $(k+1)$-th expression.
\item \textbf{Isolate the Hypothesis:} Regroup the expanded terms to separate out the $k$-th expression.
\item \textbf{Substitute and Factor:} Replace the $k$-th expression with $d \cdot m$, and factor $d$ out of the remaining terms.
\end{enumerate}

\begin{theorem}
For all integers $n \ge 1$, $3$ divides $n^3 + 2n$.
\end{theorem}

\begin{proof}
Let $P(n)$ be the predicate that $3 \mid (n^3 + 2n)$.

\textbf{Base Case ($n = 1$):}
For $n = 1$, the expression evaluates to:
\[
1^3 + 2(1) = 1 + 2 = 3 = 3 \cdot 1.
\]
Because $3$ is an integer multiple of $3$, the base case $P(1)$ is verified.

\textbf{Inductive Step:}
Let $k \ge 1$ be an arbitrary integer. Assume $P(k)$ is true: there exists an integer $m$ such that
\[
k^3 + 2k = 3m. \tag{Induction Hypothesis}
\]
We must prove $P(k+1)$: that $(k+1)^3 + 2(k+1)$ is an integer multiple of $3$.

We expand $(k+1)^3 + 2(k+1)$:
\begin{align*}
(k+1)^3 + 2(k+1) 
&= (k^3 + 3k^2 + 3k + 1) + (2k + 2) \\
&= k^3 + 3k^2 + 5k + 3.
\end{align*}
We isolate the expression from our induction hypothesis, $k^3 + 2k$. Grouping $k^3$ together with $2k$ leaves the remaining terms together:
\[
(k+1)^3 + 2(k+1) = (k^3 + 2k) + (3k^2 + 3k + 3).
\]
By our induction hypothesis, $k^3 + 2k = 3m$. Substituting this into our expression:
\begin{align*}
(k+1)^3 + 2(k+1) 
&= 3m + 3(k^2 + k + 1) \\
&= 3\bigl(m + k^2 + k + 1\bigr).
\end{align*}
Because $m$ and $k$ are integers, $q = m + k^2 + k + 1$ is an integer. Thus $(k+1)^3 + 2(k+1) = 3q$, proving that $3 \mid ((k+1)^3 + 2(k+1))$.

By the Principle of Mathematical Induction, $3 \mid (n^3 + 2n)$ for all integers $n \ge 1$.
\end{proof}

\begin{trapbox}
\textbf{Common Mistakes in Induction Proofs:}
\begin{enumerate}
\item \emph{Omitting or mishandling the base case:} Without a verified base case, an inductive implication proves nothing. For instance, the implication $n = n + 1 \implies n+1 = (n+1) + 1$ is a valid logical conditional, but $n = n+1$ is false for every integer because no base case holds to begin the chain.
\item \emph{Assuming the target $P(k+1)$:} You must adopt the predecessor case $P(k)$ as your working hypothesis and deduce $P(k+1)$. Assuming $P(k+1)$ at the outset is circular logic.
\item \emph{Proving the inductive step for only a specific number:} Proving $P(1) \implies P(2)$ and $P(2) \implies P(3)$ does not prove the theorem. The inductive step must be established for an \emph{arbitrary} integer $k \ge n_0$.
\item \emph{Hiding or ignoring the induction hypothesis:} If your proof of $P(k+1)$ never uses $P(k)$, you are either performing direct algebra (which did not need induction in the first place) or your deduction contains a hidden hole. Always explicitly flag where the induction hypothesis is invoked.
\end{enumerate}

\textbf{P\'olya's Horse Fallacy: The Hidden Disconnect}

To see how a subtle defect in the inductive step can prove absurdities, consider George P\'olya's famous paradox: ``Prove that all horses have the same color.''
\begin{itemize}
\item \emph{Base Case ($n=1$):} In any set consisting of a single horse, every horse trivially shares the same color, because there is only one horse.
\item \emph{Inductive Step:} Assume that in any set of $k$ horses, all horses have the identical color. Now consider any set of $k+1$ horses, numbered $\{h_1, h_2, \dots, h_{k+1}\}$. 
Remove horse $h_{k+1}$; the remaining $k$ horses $\{h_1, h_2, \dots, h_k\}$ have the same color by our induction hypothesis. 
Now replace $h_{k+1}$ and remove $h_1$; the remaining $k$ horses $\{h_2, \dots, h_{k+1}\}$ also have the same color. 
Because both groups overlap, all $k+1$ horses must have the same color!
\end{itemize}

Where does this deduction collapse? Look at the transition from $k=1$ to $k+1=2$ illustrated in Figure~17.3!
\end{trapbox}

\begin{center}
\begin{tikzpicture}[scale=0.9, >=Latex]
  % Stage box
  \draw[rounded corners=4pt, fill=gray!3, draw=gray!25] (-0.5, -0.9) rectangle (13.3, 3.4);

  % Panel 1: k=2 (3 horses) -> Overlap works
  \begin{scope}[xshift=0.3cm, yshift=0.1cm]
    \node[anchor=west, font=\footnotesize\bfseries, text=targetgreen!80!black] at (-0.2, 2.8) {1. Valid Overlap ($k=2 \implies 3$)};

    % Horse nodes
    \node[circle, draw=sourceblue, fill=sourceblue!20, thick, inner sep=2.5pt, font=\scriptsize\bfseries] (h1) at (0.8, 1.4) {$h_1$};
    \node[circle, draw=structurepurple, fill=structurepurple!25, thick, inner sep=2.5pt, font=\scriptsize\bfseries] (h2) at (2.4, 1.4) {$h_2$};
    \node[circle, draw=quotientorange, fill=quotientorange!20, thick, inner sep=2.5pt, font=\scriptsize\bfseries] (h3) at (4.0, 1.4) {$h_3$};

    % Group A oval enclosing h1, h2
    \draw[thick, sourceblue, dashed, rounded corners=8pt] (0.2, 0.9) rectangle (3.0, 1.9);
    \node[sourceblue!90!black, font=\tiny\bfseries, anchor=south] at (1.0, 1.95) {Group $A = \{h_1, h_2\}$};

    % Group B oval enclosing h2, h3
    \draw[thick, quotientorange, dashed, rounded corners=8pt] (1.8, 0.75) rectangle (4.6, 1.75);
    \node[quotientorange!90!black, font=\tiny\bfseries, anchor=north] at (3.8, 0.7) {Group $B = \{h_2, h_3\}$};

    % Overlap annotation
    \draw[->, very thick, structurepurple] (2.4, 0.1) -- (h2);
    \node[below, font=\scriptsize\bfseries, text=structurepurple!90!black, align=center] at (2.4, 0.05) {Overlap $A \cap B = \{h_2\}$ (bridge exists)\\[2pt]\footnotesize Color propagates between groups};
  \end{scope}

  % Divider
  \draw[dashed, gray!30] (6.5, -0.6) -- (6.5, 3.1);

  % Panel 2: k=1 (2 horses) -> Overlap collapses!
  \begin{scope}[xshift=7.0cm, yshift=0.1cm]
    \node[anchor=west, font=\footnotesize\bfseries, text=trapred!90!black] at (-0.2, 2.8) {2. The Fatal Gap ($k=1 \implies 2$ fails)};

    % Horse nodes
    \node[circle, draw=sourceblue, fill=sourceblue!20, thick, inner sep=2.5pt, font=\scriptsize\bfseries] (x1) at (0.8, 1.4) {$h_1$};
    \node[circle, draw=quotientorange, fill=quotientorange!20, thick, inner sep=2.5pt, font=\scriptsize\bfseries] (x2) at (4.4, 1.4) {$h_2$};

    % Group A oval enclosing only h1
    \draw[thick, sourceblue, dashed, rounded corners=8pt] (0.1, 0.8) rectangle (1.5, 2.0);
    \node[sourceblue!90!black, font=\tiny\bfseries, anchor=south] at (0.8, 2.05) {Group $A = \{h_1\}$};

    % Group B oval enclosing only h2
    \draw[thick, quotientorange, dashed, rounded corners=8pt] (3.7, 0.8) rectangle (5.1, 2.0);
    \node[quotientorange!90!black, font=\tiny\bfseries, anchor=south] at (4.4, 2.05) {Group $B = \{h_2\}$};

    % Gap indication: text and arrow well-separated
    \node[font=\footnotesize\bfseries, text=trapred] at (2.6, 1.6) {No overlap!};
    \draw[<->, thick, trapred] (1.8, 1.1) -- (3.4, 1.1);
    \node[below, font=\scriptsize\bfseries, text=trapred!90!black, align=center] at (2.6, 0.05) {Overlap $A \cap B = \emptyset$ (no bridge)\\[2pt]\footnotesize Step $1 \implies 2$ cannot fire!};
  \end{scope}
\end{tikzpicture}\\[0.4em]
{\small\textbf{Figure 17.3:} Why P\'olya's horse fallacy collapses. \emph{Left:} For $k \ge 2$, any two sub-collections of size $k$ share at least one common horse ($h_2$), which transmits color between groups. \emph{Right:} For $k=1$, the subsets $\{h_1\}$ and $\{h_2\}$ are completely disjoint. Without an overlapping bridge horse, the induction step $P(1) \implies P(2)$ cannot fire.}
\end{center}

When $k = 1$, the first group is $\{h_1\}$ and the second group is $\{h_2\}$. Their intersection is empty:
\[
\{h_1\} \cap \{h_2\} = \emptyset.
\]
There is no shared horse between the two groups when $k=1$. The inductive step $P(1) \implies P(2)$ is false. Because the implication fails at $k=1$, the principle of induction cannot be applied.

\subsection{The Reach and Limits of Induction}

Mathematical induction allows proving statements across all natural numbers:
\begin{itemize}
\item \textbf{Infinite statements in two steps:} We can verify summation formulas, divisibility laws, and algebraic inequalities across all natural numbers using a concise chain of deductions.
\item \textbf{Algorithm correctness via loop invariants:} In computer science, an induction proof on loop counter $k$ proves that a program preserves its safety and correctness conditions at every iteration.
\item \textbf{Asymptotic scaling:} In analysis, induction establishes foundational bounds (such as $2^n > n$), providing the basis for convergence tests in calculus.
\end{itemize}

Yet ordinary mathematical induction has a limitation: it is strictly \textbf{single-step and local}.

Ordinary induction demands that an object of size $k+1$ is constructed by adding a single unit onto an object of size $k$. But what happens when an object does not reduce to its immediate predecessor?
\begin{itemize}
\item \textbf{Multiplicative Factorization:} When factoring $12 = 3 \times 4$, neither $3$ nor $4$ is $11$. Knowing that $11$ can be factored into primes gives no information about how to factor $12$.
\item \textbf{Divide-and-Conquer Algorithms:} In algorithms such as merge sort or binary search, an input of size $n$ splits into subproblems of size $n/2$, not size $n-1$.
\item \textbf{Recursive Sequences:} In the Fibonacci sequence $F_{n+1} = F_n + F_{n-1}$, calculating the next value requires the two previous values, not just one.
\item \textbf{Geometric Triangulation:} Cutting an arbitrary polygon of $k+1$ vertices along an interior diagonal splits the shape into two smaller polygons, both strictly smaller than $k+1$, but neither guaranteed to be of size $k$.
\end{itemize}

When subproblems jump backward by more than one step, ordinary induction stalls. We need an induction principle that assumes all previous cases $P(1), P(2), \dots, P(k)$.

In Chapter~18, we introduce \textbf{Strong Induction and Recursive Thinking}, showing how assuming the statement for all smaller values allows proving the Fundamental Theorem of Arithmetic and analyzing recursive algorithms.

\subsection{Exercises}

\begin{enumerate}
\item \textbf{Mental Check (Base Case Location):} Test the inequality $2^n > n^2$ for small positive integers:
\begin{enumerate}[label=(\alph*)]
\item Evaluate both sides for $n = 1, 2, 3, 4, 5, 6$. For which values does the inequality hold, and for which does it fail?
\item Identify the smallest integer $n_0$ where the inequality holds permanently.
\item Prove by induction that $2^n > n^2$ for all integers $n \ge 5$. (Where in your inductive step do you use the fact that $k \ge 5$?)
\end{enumerate}

\item \textbf{Sum of Odds:} Prove by induction that the sum of the first $n$ odd positive integers equals $n^2$:
\[
1 + 3 + 5 + \dots + (2n - 1) = n^2 \quad \text{for all integers } n \ge 1.
\]
(Draw a $k \times k$ square of dots and see how adding an L-shaped outer border of $2k + 1$ dots expands it into a $(k+1) \times (k+1)$ square.)

\item \textbf{Divisibility Practice:} Prove by induction that $3 \mid (4^n - 1)$ for all integers $n \ge 1$. (Follow the four-step strategy: state the hypothesis $4^k - 1 = 3m$, write $4^{k+1} - 1 = 4 \cdot 4^k - 1$, and hunt for the expression $4^k - 1$.)

\item \textbf{Factorial Growth:} Prove by induction that $n! > 2^n$ for all integers $n \ge 4$.

\item \textbf{Bernoulli's Inequality:} Let $x$ be a real number with $x > -1$ and $x \neq 0$. Prove by induction that for all integers $n \ge 2$,
\[
(1 + x)^n > 1 + nx.
\]
Where in your proof do you rely on the assumption that $1 + x > 0$?

\item \textbf{Geometric Challenge (L-Tromino Tiling):} An L-tromino is a tile made of three $1 \times 1$ squares joined in an L-shape. A $2^n \times 2^n$ grid of squares has exactly one $1 \times 1$ square removed at an arbitrary position. Prove by induction on $n \ge 1$ that the remaining $2^{2n} - 1$ squares can always be completely covered without overlap by L-trominoes.
\emph{Hint:} Divide the $2^{k+1} \times 2^{k+1}$ grid into four $2^k \times 2^k$ quadrants. Exactly one quadrant contains the missing square. Place a single L-tromino at the central intersection of the other three quadrants so that it covers one corner of each!
\end{enumerate}

\newpage
